\paragraph{6. Errorless Average-Case Hardness of NP.}
In this section, we show the hardness part of our oracle separation. First, we show the existence of AIPRG relative to $\mathcal O_a$. Then we show the other hardness results, including the errorless average-case hardness for $\mathsf{NP}$, as corollaries.
We use the following theorem, which shows the existence of PRGs based on a random oracle.
\paragraph{Theorem 18 ([45]).}
For each $n\in\mathbb N$, let $R_n:\{0,1\}^n\to\{0,1\}^n$ be a random function oracle, i.e., $R_n$ is selected uniformly at random from $\{f:\{0,1\}^n\to\{0,1\}^n\}$.
There exist a polynomial-time deterministic oracle machine $G^?$ and constants $c\geq1$ and $b,\epsilon>0$ satisfying the following: For any $n\in\mathbb N$ and $x\in\{0,1\}^{cn}$, $G^{R_n}(x)$ generates a binary string of length $4cn$, and all oracle circuits $C^?$ of size $2^{bn}$ satisfy that
\[\left|\Pr_{U_{cn}}[C^{R_n}(G^{R_n}(U_{cn}))=1]-\Pr_{U_{4cn}}[C^{R_n}(U_{4cn})=1]\right|\leq2^{-\epsilon n},\]
with probability at least $1-2^n$ over the choice of $R_n$.
Furthermore, the result above is relativized, i.e., the above holds in the presence of an arbitrary oracle $O$ independent of the choice of $R_n$.
Now, we show the existence of AIPRG relative to $\mathcal O_a$.
\paragraph{Theorem 19.}
Let $a>0$ be an arbitrary constant. With probability 1 over the choice of $\mathcal O_a$, there exists an AIPRG $G^{\mathcal O_a}=\{G_z^{\mathcal O_a}\}_{z\in\{0,1\}^*}$ against $\mathsf{SIZE}^{\mathcal O_a}[2^{\epsilon an/\log n}]$ for some absolute constant $\epsilon>0$, where $G_z^{\mathcal O_a}:\{0,1\}^{|z|}\to\{0,1\}^{3|z|}$ for each $z\in\{0,1\}^*$.
\begin{proof}
Let $G^?$ and $c$ be the oracle machine and the constant in Theorem 18, respectively. Then, we define an AIPRG $G'=\{G'_z\}_{z\in\{0,1\}^*}$ as $G_z'^{\mathcal O}(x)=G^{F_{z'}}(x')_{\leq3|z|}$, where $x\in\{0,1\}^{|z|}$, $z'=z_{\leq\lfloor|z|/c\rfloor}$, $x'=x_{\leq c|z'|}$, and $F_{z'}:\{0,1\}^{|z'|}\to\{0,1\}^{|z'|}$ is defined as $F_{z'}(y)=\mathcal F(z',y,1)\circ\cdots\circ\mathcal F(z',y,|z'|)$. The validity of the truncation is verified as that $|x'|=c|z'|\leq c\cdot|z|/c=|z|=|x|$ and $|G^{F_{z'}}(x')|=4|x'|=4c|z'|\geq4c(|z|/c-1)\geq4|z|-4c\geq3|z|$ for any $z$ with $|z|\geq4c$.
Let $\epsilon=1/2c$ and $s(n)=2^{\epsilon an/\log n}$. We show that $G'$ above is an AIPRG against $\mathsf{SIZE}^{\mathcal O}[s(n)]$. Suppose there exists a family $C=\{C_n\}_{n\in\mathbb N}$ of oracle circuits of size $s(n)$ that breaks $G'$, i.e., for any sufficiently large $n\in\mathbb N$ and any $z\in\{0,1\}^n$,
\[\left|\Pr_{U_n}[C_n^{\mathcal O}(z,G_z'^{\mathcal O}(U_n))=1]-\Pr_{U_{3n}}[C_n^{\mathcal O}(z,U_{3n})=1]\right|>\frac1{\operatorname{poly}(n)}.\]
Fix $n\in\mathbb N$ arbitrarily, and let $n'=\lfloor n/c\rfloor$. We consider an arbitrary choice of $\mathcal O$ except for the values of $\rho_{n',i_{\max}(n')}$ (we write this condition as $R$ for convenience). Fix $z\in\{0,1\}^n$ such that $z'=z_{\leq n'}\in S_{n',i_{\max}(n')-1}$ arbitrary (where $S_{\cdot,\cdot}$ is the random set in the oracle construction). We refer to such $z$ as a hard index.
Since the size of $C_n$ is at most $s(n)=2^{an/(2c\log n)}\leq2^{a\lfloor n/c\rfloor/\log\lfloor n/c\rfloor}=t(n')$ for sufficiently large $n$ (where $t$ is the time-bound function in the oracle construction), the answers to queries by $C_n$ of the form $\mathcal A(M,y,1^{T^{2^c}})$ do not depend on the values of $F_{z'}$ and are determined by condition $R$ because they are determined only by the restrictions $\rho_{\cdot,j}$ for
\[j\leq c^{-1}\log\log T\leq c^{-1}\log\log s(n)^{1/2^c}\leq c\log\log t(n')-1<i_{\max}(n').\]
Now, we consider an arbitrary choice of $\rho_{n',i_{\max}(n')}$ except for the values of $F_{z'}$ and denote this condition by $R'$. It is not hard to verify that $F_{z'}$ is selected uniformly at random even under the conditions $R$ and $R'$. Therefore, under the conditions $R$ and $R'$, we can identify the query access to $\mathcal O$ by $C$ with the query access to another oracle $O'$ (determined only by $R$ and $R'$) and a random function oracle $F_{z'}$ that are selected independently of $O'$.
For any $n\in\mathbb N$, let $E_n$ be an event (over the choice of $F_{z'}$) that there exists a circuit $C'$ of size $2^{bn'}$, where $b$ represents the constant in Theorem 18, such that
\begin{align*}
&\left|\Pr_{U_n}[C'^{\mathcal O}(G_z'^{F_{z'}}(U_n))=1]-\Pr_{U_{3n}}[C'^{\mathcal O}(U_{3n})=1]\right|\\
&=\left|\Pr_{U_n}[C'^{O',F_{z'}}(G^{F_{z'}}(U_n))=1]-\Pr_{U_{3n}}[C'^{O',F_{z'}}(U_{3n})=1]\right|>\frac1{\operatorname{poly}(n)}.
\end{align*}
Then, by Theorem 18 (relative to $O'$), we have $\Pr_{\mathcal O}[E_n\mid R,R']\leq2^{-\Omega(n)}$. By the Borel--Cantelli lemma, $E_n$ occurs only for finitely many $n$'s with probability 1 over the choice of $\mathcal O$ conditioned on $R,R'$ (i.e., the choice of $F_{z'}$). By taking the expectation over $R,R'$, we can show that, with probability 1 over the choice of $\mathcal O$, there is no family $C'$ of $2^{bn'}$-size circuits satisfying that for any sufficiently large $n\in\mathbb N$, there exists a hard index $z\in\{0,1\}^n$ such that
\[\left|\Pr_{U_n}[C'^{\mathcal O}(G_z'^{F_{z'}}(U_n))=1]-\Pr_{U_{3n}}[C'^{\mathcal O}(U_{3n})=1]\right|>\frac1{\operatorname{poly}(n)}.\]
However, the circuit $C$ in the assumption violates this statement by embedding a hard index $z$ as auxiliary-input because the size is at most $n+s(n)=O(2^{(a/2c)\cdot n/\log n})=o(2^{bn'})$. Therefore, we conclude that, with probability 1 over the choice of $\mathcal O$, there is no such a circuit $C$ of size $s(n)$, and $G'$ is an AIPRG against $\mathsf{SIZE}[s(n)]$.
\end{proof}
Next, AIPRG implies HSG by regarding the auxiliary-input as a part of the hidden input to HSG. It is not hard to verify the security (refer to [34, Lemma 18] for the formal proof).
\paragraph{Corollary 21.}
Let $a>0$ be an arbitrary constant. With probability 1 over the choice of $\mathcal O_a$, there exists an HSG $G^{\mathcal O_a}=\{G_n^{\mathcal O_a}\}_{n\in\mathbb N}$ against $\mathsf{SIZE}^{\mathcal O_a}[2^{\epsilon an/\log n}]$ for some absolute constant $\epsilon>0$, where $G_n^{\mathcal O_a}:\{0,1\}^{2n}\to\{0,1\}^{3n}$ for each $n\in\mathbb N$.
Furthermore, the existence of HSGs implies the errorless average-case hardness of $\mathsf{NP}$, as observed in [21] in the context of natural proofs and the average-case complexity of the minimum circuit size problem.
\paragraph{Corollary 22.}
Let $a>0$ be an arbitrary constant. With probability 1 over the choice of $\mathcal O_a$, $\mathsf{DistNP}^{\mathcal O_a}\not\subseteq\mathsf{AvgSIZE}^{\mathcal O_a}[2^{\epsilon an/\log n}]$ for some absolute constant $\epsilon>0$.
\begin{proof}
Let $G^{\mathcal O}$ and $\epsilon$ be the HSG and the constant in Corollary 21, respectively. Then, we define the language $L^{\mathcal O}$ as $L^{\mathcal O}:=\operatorname{Im}G^{\mathcal O}$. Obviously, $L^{\mathcal O}\in\mathsf{NP}^{\mathcal O}$ and $(L^{\mathcal O},\{U_n\}_{n\in\mathbb N})\in\mathsf{DistNP}^{\mathcal O}$. Thus, it is sufficient to show that $(L^{\mathcal O},\{U_n\}_{n\in\mathbb N})\notin\mathsf{Avg}_{1/4}\mathsf{SIZE}^{\mathcal O}[2^{\epsilon an/(4\log n)}]$.
For contradiction, we assume that $(L^{\mathcal O},\{U_n\}_{n\in\mathbb N})\in\mathsf{Avg}_{1/4}\mathsf{SIZE}^{\mathcal O}[2^{\epsilon an/(4\log n)}]$. Then, there exists a family $C=\{C_n\}_{n\in\mathbb N}$ of $O(2^{\epsilon an/(4\log n)})$-size oracle circuits for $(L^{\mathcal O},\{U_n\}_{n\in\mathbb N})$, i.e., for any sufficiently large $n\in\mathbb N$,
\[\Pr_{y\leftarrow_u\{0,1\}^{3n}}[C_{3n}^{\mathcal O}(y)=\bot]\leq1/4\]
and for each $y\in\{0,1\}^{3n}$,
\[C_{3n}^{\mathcal O}(y)\in\{\mathbf1(y\in L^{\mathcal O}),\bot\}.\]
Note that the size of $C_{3n}$ is at most $O(2^{\epsilon a\cdot3n/(4\log3n)})$.
Next, we define an adversary $C'$ for $G^{\mathcal O}$ as follows: for a given $y\in\{0,1\}^{3n}$ (i.e., the length of seed is $2n$), $C_{2n}'^{\mathcal O}$ simulates $C_{3n}^{\mathcal O}(y)$, and if $C_{3n}$ returns 1 or $\bot$, then $C'_{2n}$ outputs 0; otherwise (i.e., if $C_{3n}$ returns 0), $C'_{2n}$ outputs 1. Then, based on the aforementioned inequalities, it is not hard to verify that for any sufficiently large $n\in\mathbb N$,
\begin{align*}
\Pr_{y\leftarrow_u\{0,1\}^{3n}}[C_{2n}'^{\mathcal O}(y)=0]
&\leq\frac14+\frac{|\{G^{\mathcal O}(x):x\in\{0,1\}^{2n}\}|}{|\{0,1\}^{3n}|}\\
&\leq\frac14+\frac{2^{2n}}{2^{3n}}<\frac12,
\end{align*}
and for any $x\in\{0,1\}^{2n}$, we have $C_{3n}^{\mathcal O}(G^{\mathcal O}(x))\in\{1,\bot\}$ and $C_{2n}'^{\mathcal O}(G^{\mathcal O}(x))=0$.
Therefore, $C'$ succeeds in avoiding $\operatorname{Im}G^{\mathcal O}$, and the size is bounded above by $O(2^{\epsilon an/\log n})=O(2^{\epsilon a(2n)/\log(2n)})$. This contradicts Corollary 21. Thus, we conclude that $(L^{\mathcal O},\{U_n\}_{n\in\mathbb N})\notin\mathsf{Avg}_{1/4}\mathsf{SIZE}^{\mathcal O}[2^{\epsilon an/(4\log n)}]$.
\end{proof}
